\section{Identified Gap}

\subsection{Limitations of Centralized Solutions}

Centralized collaboration platforms provide strong support for
membership management, synchronization, file compatibility,
and usability. However, these advantages are achieved through
reliance on centralized infrastructure for storage,
coordination, and access management.

As a result, collaboration depends on the availability of a
central service and its associated network connectivity.
Although these platforms support basic version recovery, they
generally provide limited facilities for long-term revision
tracking and conflict management when compared with dedicated
version control systems.

Furthermore, ownership and control of shared data remain tied
to the centralized platform. This model may be unsuitable for
temporary collaborative groups seeking lightweight and
serverless file sharing.

\subsection{Limitations of Collaborative Version Control Systems}

Collaborative version control systems provide strong support
for revision tracking, rollback, and conflict resolution.
These capabilities make them highly effective for managing the
evolution of shared project artifacts over time.

However, such systems are primarily designed for controlled
revision management rather than transparent file
synchronization. Changes are typically exchanged through
explicit operations such as commits, pushes, pulls, or
updates, requiring users to actively participate in the
synchronization process.

In addition, version control workflows often introduce a
higher level of complexity than ordinary file sharing
solutions. Repository management, branching, merging, and
conflict resolution require specialized tools and technical
knowledge that may be unnecessary for small temporary
collaboration groups.

Furthermore, support for large binary files and multimedia
assets may be limited depending on the system being used. As a
result, collaborative version control systems are well suited
for managing project history but are less suitable as
lightweight directory synchronization solutions.

\subsection{Limitations of Existing Peer-to-Peer Solutions}

Existing peer-to-peer approaches address different aspects of
collaborative file management, but none provide comprehensive
support across all evaluation dimensions.

Collaborative version control systems such as
\textit{Git}, \textit{SVN}, and \textit{Perforce} offer strong
history management through version tracking, rollback, and
conflict resolution. However, they rely on explicit workflows
and are not designed for transparent directory
synchronization.

Conversely, peer-to-peer synchronization systems such as
\textit{Syncthing} and \textit{Resilio Sync} provide automatic
directory synchronization and strong filesystem integration.
Nevertheless, their history management capabilities are
typically limited to basic version retention and recovery
mechanisms.

Local sharing technologies such as \textit{AirDrop} and
\textit{Nearby Share} further simplify file exchange, but they
do not provide persistent synchronization, collaborative
history management, or long-term shared workspaces.

\subsection{Limitations of Local Sharing Technologies}

Local sharing technologies simplify file exchange by enabling
direct device-to-device transfers between nearby users.
Systems such as \textit{AirDrop} and \textit{Nearby Share}
provide convenient mechanisms for sharing files without
requiring centralized infrastructure or complex setup.

However, these technologies are designed primarily for
one-time transfers rather than persistent collaboration. They
do not maintain synchronized directories, automatically
propagate subsequent modifications, or preserve revision
history.

In addition, sharing relationships are typically temporary and
must be re-established whenever files are exchanged. As a
result, local sharing technologies are highly effective for
quick file transfers but do not provide the functionality
required for ongoing collaborative file management.

\subsection{The Gap}

The limitations discussed in the previous sections
show that existing approaches provide different strengths for
collaborative file management, but none fully satisfies the
combination of requirements. 

Table~\ref{tab:gap-summary} summarizes the performance of the
reviewed approaches across the five evaluation dimensions
introduced in Chapter~3: membership management,
synchronization, history management, file support, and
filesystem transparency and usability. The comparison shows
that each category focuses on a particular subset of these
dimensions.

Centralized collaboration platforms provide strong support for
membership management, synchronization, file support, and
usability, but offer only limited history management.
Collaborative version control systems provide comprehensive
history management, but provide weaker support for
synchronization, file support, and transparent filesystem
integration. Peer-to-peer synchronization systems provide
strong synchronization and usability, but only limited support
for membership management and version recovery. Local sharing
technologies simplify file exchange and provide strong
usability, but do not support persistent synchronization or
history management.

Consequently, the gap identified in this study is not the
absence of any individual capability. Rather, the gap is the
absence of a lightweight solution that simultaneously provides
strong support across all five evaluation dimensions.
Existing systems typically optimize one aspect of
collaborative file management at the expense of others.

Therefore, there remains a need for a lightweight peer-to-peer
directory synchronization system that combines decentralized
collaboration, automatic synchronization, basic version
recovery, support for heterogeneous project files, and
transparent filesystem integration within a single solution.

Local Sharing Technologies
\begin{table}[htbp]
\centering
\footnotesize
\caption{Gap Summary Across Evaluation Dimensions}
\label{tab:gap-summary}

\begin{tabular}{
  C{2.8cm}
  C{2.5cm}
  C{2.8cm}
  C{2.5cm}
  C{2.5cm}
  C{2.5cm}
}
\toprule

\textbf{Approach}
&
\makecell{\textbf{Membership}\\\textbf{Management}}
&
\makecell{\textbf{Synchronization}}
&
\makecell{\textbf{History}\\\textbf{Management}}
&
\makecell{\textbf{File}\\\textbf{Support}}
&
\makecell{\textbf{Filesystem}\\\textbf{Transparency}\\\textbf{\& Usability}}

\\

  \midrule

  Centralized Platforms
  & \high
  & \high
  & \limited
  & \high
  & \high
  \\

  Version Control Systems
  & \partialsupport
  & \limited
  & \high
  & \partialsupport
  & \limited
  \\

  P2P Synchronization Systems
  & \partialsupport
  & \high
  & \limited
  & \high
  & \high
  \\

& \limited
& \none
& \none
& \high
& \high
\\

\midrule

\textbf{Desired Solution}
& \high
& \high
& \high
& \high
& \high
\\

\bottomrule
\end{tabular}
\end{table}
